\section{Famous ``Dream Discoveries'': Fact vs\. Myth}\label{sec:famous_eurekas}

The journey from the psychoanalyst's couch to the neuroscience lab clarified that ``dreaming'' is not a single state but a family of experiences tied to specific neural architectures \citep{revonsuo2013important}. The chapters in Part~I review how cultural intuitions about dreams foreshadowed this realization, while Part~II shows how deprivation studies make clear that sleep loss is a catastrophic failure of an essential offline computation \citep{wilson1994reactivation, waters2018deprivation}. Before that science took shape, public imagination latched onto a handful of charming origin stories in which epochal discoveries supposedly arrived wholesale in a dream. They are worth recounting---and debunking---because they highlight how much explanatory weight we placed on nocturnal revelation.

\paragraph{Kekul\'e (benzene).}
In an 1890 commemorative address (the “Benzolfest”), Kekul\'e himself described an image of a serpent biting its tail that suggested cyclic structure; modern historians emphasize this was a retrospective daydream anecdote rather than a laboratory epiphany, and caution against over-literal readings \citep{benfey1958, rocke2010}.

\paragraph{Mendeleev (periodic table).}
The familiar story that Mendeleev \emph{dreamed} the periodic system is not supported by the archival record; the dream motif appears to be a later embellishment---he had long been iterating toward the table via notebooks and card-sorting schemes \citep{gordin2004}.

\paragraph{Elias Howe (sewing machine).}
The tale of cannibals brandishing spears with holes near the tip (inspiring the needle’s eye placement) is a durable piece of popular history; it traces to early 20th-century retellings rather than contemporaneous primary documentation \citep{kaempffert1924}. It remains illustrative, but should be presented explicitly as anecdote.

\paragraph{Ramanujan (mathematics)}
Ramanujan repeatedly attributed results to goddess Namagiri and reported nocturnal inspiration; reputable biographies treat this as part of his lived creative practice, while avoiding over-specific dramatization \citep{kanigel1991}.
